\documentclass[10pt,a4paper]{amsart}
\usepackage{amsmath,amssymb,amsthm,graphicx}
\usepackage{verbatim}
\newenvironment{reference}{\comment}{\endcomment}
%\newenvironment{reference}{}{}
\newenvironment{slogan}{\comment}{\endcomment}
\newenvironment{history}{\comment}{\endcomment}

% For commutative diagrams we use Xy-pic
\usepackage[all]{xy}

% We use 2cell for 2-commutative diagrams.
\xyoption{2cell}
% Theorem environments.
%
\theoremstyle{plain}
\newtheorem{theorem}[subsection]{Theorem}
\newtheorem{proposition}[subsection]{Proposition}
\newtheorem{lemma}[subsection]{Lemma}

\theoremstyle{definition}
\newtheorem{definition}[subsection]{Definition}
\newtheorem{example}[subsection]{Example}
\newtheorem{exercise}[subsection]{Exercise}
\newtheorem{situation}[subsection]{Situation}

\theoremstyle{remark}
\newtheorem{remark}[subsection]{Remark}
\newtheorem{remarks}[subsection]{Remarks}

\numberwithin{equation}{subsection}

% Macros
%
\def\lim{\mathop{\mathrm{lim}}\nolimits}
\def\colim{\mathop{\mathrm{colim}}\nolimits}
\def\Spec{\mathop{\mathrm{Spec}}}
\def\Hom{\mathop{\mathrm{Hom}}\nolimits}
\def\Ext{\mathop{\mathrm{Ext}}\nolimits}
\def\SheafHom{\mathop{\mathcal{H}\!\mathit{om}}\nolimits}
\def\SheafExt{\mathop{\mathcal{E}\!\mathit{xt}}\nolimits}
\def\Sch{\mathit{Sch}}
\def\Mor{\mathop{\mathrm{Mor}}\nolimits}
\def\Ob{\mathop{\mathrm{Ob}}\nolimits}
\def\Sh{\mathop{\mathit{Sh}}\nolimits}
\def\NL{\mathop{N\!L}\nolimits}
\def\CH{\mathop{\mathrm{CH}}\nolimits}
\def\proetale{{pro\text{-}\acute{e}tale}}
\def\etale{{\acute{e}tale}}
\def\QCoh{\mathit{QCoh}}
\def\Ker{\mathop{\mathrm{Ker}}}
\def\Im{\mathop{\mathrm{Im}}}
\def\Coker{\mathop{\mathrm{Coker}}}
\def\Coim{\mathop{\mathrm{Coim}}}

% Boxtimes
%
\DeclareMathSymbol{\boxtimes}{\mathbin}{AMSa}{"02}

%
% Macros for moduli stacks/spaces
%
\def\QCohstack{\mathcal{QC}\!\mathit{oh}}
\def\Cohstack{\mathcal{C}\!\mathit{oh}}
\def\Spacesstack{\mathcal{S}\!\mathit{paces}}
\def\Quotfunctor{\mathrm{Quot}}
\def\Hilbfunctor{\mathrm{Hilb}}
\def\Curvesstack{\mathcal{C}\!\mathit{urves}}
\def\Polarizedstack{\mathcal{P}\!\mathit{olarized}}
\def\Complexesstack{\mathcal{C}\!\mathit{omplexes}}
% \Pic is the operator that assigns to X its picard group, usage \Pic(X)
% \Picardstack_{X/B} denotes the Picard stack of X over B
% \Picardfunctor_{X/B} denotes the Picard functor of X over B
\def\Pic{\mathop{\mathrm{Pic}}\nolimits}
\def\Picardstack{\mathcal{P}\!\mathit{ic}}
\def\Picardfunctor{\mathrm{Pic}}
\def\Deformationcategory{\mathcal{D}\!\mathit{ef}}
\usepackage{cleveref}
\crefname{originalreference}{original reference}{original references}
\Crefname{originalreference}{Original reference}{Original references}
\makeatletter
\newlabel{situation-iso}{{05MH}{source}{}{sourcepointer-1}{}}
\newlabel{situation-iso@cref}{{[originalreference][0][]05MH}{[1][0][]source}{}{}{}}
\makeatother
\crefname{originallemma}{original lemma}{original lemmas}
\Crefname{originallemma}{Original lemma}{Original lemmas}
\makeatletter
\newlabel{spaces-morphisms-lemma-support-finite-type}{{07TZ}{source}{}{sourcepointer-2}{}}
\newlabel{spaces-morphisms-lemma-support-finite-type@cref}{{[originallemma][0][]07TZ}{[1][0][]source}{}{}{}}
\makeatother
\crefname{originalsection}{original section}{original sections}
\Crefname{originalsection}{Original section}{Original sections}
\makeatletter
\newlabel{spaces-morphisms-section-flat}{{03MK}{source}{}{sourcepointer-3}{}}
\newlabel{spaces-morphisms-section-flat@cref}{{[originalsection][0][]03MK}{[1][0][]source}{}{}{}}
\makeatother
\title{Functor equality for the iso, injective, surjective and zero functors under flat base change when the image contains the supports}
\author{The Stacks Project Authors}
\date{stacks-project:spaces-flat.tex:local-snapshot-20261012; GNU FDL 1.2 or later}
\begin{document}
\maketitle
\noindent\textit{Editorial note.} All delivered mathematical text is the author's own, copied from the pinned Stacks Project chapter file spaces-flat.tex; no statement or proof step was rewritten, reordered or removed.\par\medskip
\noindent\textit{Editorial note.} The counted proof is the author's complete argument: a general faithful flatness argument for pullback stability of the four properties, a support-localized test criterion, and commutation of supports with base change under the finite type hypotheses\par\medskip
\noindent\textit{Editorial note.} The label lemma-iso-go-up names the counted statement, and the proof environment immediately following it in the source is delivered verbatim as its proof.\par\medskip
\noindent\textit{Editorial note.} The counted statement is carried by the author's own lemma environment in the chapter source.\par\medskip
\noindent\textit{Allowed standard prerequisites (editorial).}\begin{itemize}
\item The situation fixing a map of quasi-coherent modules u: F -> G on an algebraic space over B and the functors F\_iso, F\_inj, F\_surj, F\_zero on Sch/B classifying points where the pulled back map is an isomorphism, injective, surjective, or zero
\item Faithful flatness of the local ring maps along a surjective flat morphism of algebraic spaces, and the resulting stability of injectivity, surjectivity, being zero, and being an isomorphism under pullback
\item Commutation of scheme theoretic supports with base change for finite type quasi-coherent modules
\end{itemize}\medskip
\ifdefined\hypertarget\hypertarget{sourcepointer-1}{}\fi
\noindent\textit{Source reference.} Original reference 05MH: sources/stacks-spaces-flat/pinned-original-source.tex, source lines 1233--1233. Prerequisite pointer invoked at proof line(s) 1292; it names a standard fact of the same chapter that the counted argument uses, and every step of the counted proof built on it is retained verbatim in the delivered body.\par\medskip
\ifdefined\hypertarget\hypertarget{sourcepointer-2}{}\fi
\noindent\textit{Source reference.} Original lemma 07TZ: sources/stacks-chow/tags.txt, source lines 9954--9954. Prerequisite pointer invoked at proof line(s) 1325; it names a standard fact of a sibling chapter that the counted argument uses, and every step of the counted proof built on it is retained verbatim in the delivered body.\par\medskip
\ifdefined\hypertarget\hypertarget{sourcepointer-3}{}\fi
\noindent\textit{Source reference.} Original section 03MK: sources/stacks-chow/tags.txt, source lines 4752--4752. Prerequisite pointer invoked at proof line(s) 1319; it names a standard fact of a sibling chapter that the counted argument uses, and every step of the counted proof built on it is retained verbatim in the delivered body.\par\medskip
\begin{lemma}
\label{lemma-iso-go-up}
In Situation \ref{situation-iso} let $X' \to X$ be a flat morphism
of algebraic spaces. Denote $u' : \mathcal{F}' \to \mathcal{G}'$
the pullback of $u$ to $X'$. Denote $F'_{iso}$, $F'_{inj}$, $F'_{surj}$,
$F'_{zero}$ the functors on $\Sch/B$ associated to $u'$.
\begin{enumerate}
\item If $\mathcal{G}$ is of finite type and the image of $|X'| \to |X|$
contains the support of $\mathcal{G}$, then $F_{surj} = F'_{surj}$
and $F_{zero} = F'_{zero}$.
\item If $\mathcal{F}$ is of finite type and the image of $|X'| \to |X|$
contains the support of $\mathcal{F}$, then $F_{inj} = F'_{inj}$
and $F_{zero} = F'_{zero}$.
\item If $\mathcal{F}$ and $\mathcal{G}$ are of finite type and the image of
$|X'| \to |X|$ contains the supports of $\mathcal{F}$ and $\mathcal{G}$,
then $F_{iso} = F'_{iso}$.
\end{enumerate}
\end{lemma}

\begin{proof}
let $v : \mathcal{H} \to \mathcal{E}$ be a map of quasi-coherent
modules on an algebraic space $Y$ and let $\varphi : Y' \to Y$ be a
surjective flat morphism of algebraic spaces, then $v$ is
an isomorphism, injective, surjective, or zero if and only if $\varphi^*v$ is
an isomorphism, injective, surjective, or zero. Namely,
for every $y \in |Y|$ there exists a $y' \in |Y'|$ and the map
of local rings
$\mathcal{O}_{Y, \overline{y}} \to \mathcal{O}_{Y', \overline{y'}}$
is faithfully flat (see
Morphisms of Spaces, Section \ref{spaces-morphisms-section-flat}).
Of course, to check for injectivity or being zero it suffices to look
at the points in the support of $\mathcal{H}$, and to check for
surjectivity it suffices to look at points in the support of $\mathcal{E}$.
Moreover, under the finite type assumptions as in the statement of
the lemma, taking the supports commutes with base change, see
Morphisms of Spaces, Lemma \ref{spaces-morphisms-lemma-support-finite-type}.
Thus the lemma is clear.
\end{proof}

\end{document}
